% generated by scripts/ledger_appendix.py
\begin{longtable}{@{}>{\raggedright\arraybackslash}p{2.7cm}>{\raggedright\arraybackslash}p{7.9cm}>{\raggedright\arraybackslash}p{1.3cm}>{\raggedright\arraybackslash}p{2.5cm}@{}}
\caption{Claims-ledger excerpt: up to two architecture-comparison claims per priority-1 core paper. Regimes: P matched parameters, D matched effective depth, IF matched inference FLOPs, TF matched training FLOPs, Dat matched data, U unmatched. Claim wording is paraphrased from the extraction records; each is backed by a verified quotation in the released ledger.}\label{tab:ledger_excerpt}\\
\toprule Paper & Claim & Regime & Scale \\ \midrule \endfirsthead
\toprule Paper & Claim & Regime & Scale \\ \midrule \endhead
\bottomrule \endfoot
ACT~\citep{ax1603.08983} & On parity, the network without ACT fails (about 40\% error) while ACT with low time penalty gets below 5\% error. & P, Dat & 128-unit RNN \\
ACT~\citep{ax1603.08983} & For character-level language modeling ACT gives little gain; ponder is concentrated at spaces and sentence ends. & P, Dat & 1500-cell LSTM \\
\citep{ax2603.08391} & Adaptive looping (Nmax=3) reduces math BPB by 22\% vs the base model and beats a 36-layer iso-FLOP model on math BPB by 6.4\% with one-third of the layers. & P, D & \textasciitilde{}200M \\
\citep{ax2603.08391} & More loops add little beyond Nmax=3 (Loop-7 improves 1.7\% over Loop-3) and commonsense slightly declines with more loops; the 36-layer iso-FLOP model remains better on commonsense. & P & \textasciitilde{}200M \\
ALBERT~\citep{ax1909.11942} & At matched depth/width (fewer unique parameters), sharing all layers hurts downstream average by 1.5 (E=128) to 2.5 (E=768) points vs not-shared. & D, Dat & ALBERT-base (12M-108M) \\
ALBERT~\citep{ax1909.11942} & With roughly matched training wall-clock time, ALBERT-xxlarge beats BERT-large by 1.5 average points despite fewer parameters (but it is \textasciitilde{}3x slower per step). & TF, U & ALBERT-xxlarge 233M vs BERT-large 334M \\
Attractor Models~\citep{ax2605.12466} & On Sudoku-Extreme/Maze-Hard, Attractor Models improve with scale while TRM collapses at 27M. & P, U & 27M \\
Attractor Models~\citep{ax2605.12466} & At matched parameters, data, optimizer and schedule, Attractor Models beat both a standard Transformer and the looped Parcae on perplexity and downstream accuracy at 140M-770M. & P, Dat & 770M \\
CoTFormer~\citep{ax2310.10845} & At equal parameters and equal number of repeats, CoTFormer beats Block Universal Transformer in perplexity, and the gap widens with more repeats (at the cost of extra attention compute). & P, D, Dat & 24-layer base x 5 repeats \\
CoTFormer~\citep{ax2310.10845} & Router-based token allocation gives a much better compute-accuracy trade-off than uniformly truncating repeats at inference; but adaptive training lags a fixed-5-repeat model at equal inference compute (23.83 vs 23.19). & IF & 24 layers, max 5 repeats \\
DEQ~\citep{ax1909.01377} & At similar parameter counts DEQs match or slightly beat Transformer-XL and TrellisNet on WikiText-103 while using 80-88\% less memory. & P, Dat & 172M total / 43M non-embedding \\
DEQ~\citep{ax1909.01377} & DEQs are slower than the corresponding explicit deep networks (runtime ratios above 1x in Table 4), though computation can be traded via tolerance or Broyden iteration limits. & P & medium \\
DT-Recall (Deep Thinking with recall)~\citep{ax2202.05826} & Trained on 9x9 mazes, DT-Recall with progressive loss solves 74\% of 201x201 mazes; feed-forward and prior DT systems completely fail on 59x59. & D, Dat & conv net \\
DT-Recall (Deep Thinking with recall)~\citep{ax2202.05826} & Recall is what prevents overthinking; recalled models appear to converge to a fixed point rather than becoming unstable as iterations grow. & Dat & conv net \\
ETD~\citep{ax2510.07358} & The best block position depends on task category: 7-4*2-5 is best for math, while an earlier block (3-4*2-9) is best for reading comprehension. & P, Dat, IF, TF & 1B \\
ETD~\citep{ax2510.07358} & Which layers to loop matters: at equal FLOPs, looping the 4 knee-selected layers beats looping all layers or the middle 12 layers. & P, IF, Dat & 1B \\
HRM~\citep{ax2506.21734} & On Sudoku-Extreme-Full, increasing width does not help while depth helps; standard (and recurrent) Transformers saturate with depth while HRM keeps improving. & D & 872M (largest width-scaled baseline) \\
HRM~\citep{ax2506.21734} & Against a same-size Transformer trained with the identical setup ('Direct pred'), swapping in HRM's hierarchy plus ACT more than doubles ARC-AGI-1 performance. & P, Dat & 27M \\
HRM-Text~\citep{ax2605.20613} & Looped Transformers and RINS beat same-size Transformers, but their advantage over a larger FLOP-matched Transformer is inconsistent. & P, TF & 1B looped vs 3B dense \\
HRM-Text~\citep{ax2605.20613} & TRM (shared H/L weights) trained at 1B with reduced recursions is unstable; at 0.6B with the same H2L3 recursion it roughly matches HRM 0.6B but needs about 2x the training FLOPs. & P, TF & 1B \\
Huginn~\citep{ax2502.05171} & Against a non-recurrent twin trained on the same 180B tokens, the recurrent model is much better on harder tasks (ARC-C, GSM8K), with similar recall tasks; parameter/FLOP matching of the twin is not stated. & Dat & 3.5B \\
Huginn~\citep{ax2502.05171} & Test-time recurrence improves reasoning benchmarks up to compute equivalent to a 50B-parameter model; comparisons are to public models trained on other data. & U & 3.5B \\
Iso-Depth phi~\citep{ax2604.21106} & At matched training compute and matched effective depth, replacing unique blocks with shared recurrences raises validation loss; one recurrence is worth r\textasciicircum{}0.46 in unique-parameter terms. & TF, D & 2.15e19 FLOPs grid; 47B-token extrapolation \\
Iso-Depth phi~\citep{ax2604.21106} & The compute-optimal loss gap of looped vs non-looped grows with r and persists at \textasciitilde{}20x the grid compute (0.061 nats at r=4, 47B tokens). & D, Dat & s=34 (d\_model 2,176), 47B tokens \\
Loop Scaling Laws~\citep{ax2609.40316} & At trillion-token scale and matched training compute, a 1.3B-total looped MoE with law-derived recurrence matches a 2.9B-total non-looped MoE on reasoning but trails on the 14-benchmark overall score, and uses more inference compute. & TF, IF, Dat & 1.3B total / 0.3B active, \textasciitilde{}3T tokens \\
Loop Scaling Laws~\citep{ax2609.40316} & Recurrence adds effective parameters only up to a finite asymptote; MoE sparsity raises that asymptote and delays saturation (routing reaches different experts across passes). & P, Dat & 1.0B active, 500B tokens \\
LoopCoder-v2~\citep{ax2606.18023} & With identical parameters, 18T-token pretraining data and SFT, the two-loop PLT improves broadly over the non-looped baseline on code generation, code reasoning, agentic SWE and tool use. & P, Dat & 7B, 18T tokens \\
LoopCoder-v2~\citep{ax2606.18023} & Loop count is strongly non-monotonic: three or more loops regress, often below the non-looped baseline (R=3 drops to 27.6\% on SWE-bench Verified). & P, Dat & 7B, 18T tokens \\
\citep{ax2608.18171} & Under matched SFT, looped models do better on compositional (multi-call, nested) tool calling, while gains on isolated single API calls are smaller and model-dependent; the Ouro comparison is against non-looped models with different pretraining. & U & 2.6B \\
\citep{ax2608.18171} & Retrofitted looped models improve several compositional categories over their parents but remain much weaker than natively pretrained looped models on deeply nested workflows. & U & 1B \\
Looped TF (ICL)~\citep{ax2311.12424} & At matched effective depth (b=12 loops of 1 layer vs 12 distinct layers), the looped model does not match the standard transformer; it needs more loops than layers (about 20 iterations of 1 layer to match 8 layers). & D & 9.48M \\
Looped TF (ICL)~\citep{ax2311.12424} & Input injection is necessary for stable extrapolation in loop count: without it, the weight-tied model diverges beyond the trained number of loops; with it, the model reaches a fixed point that persists past training iterations. & P & 0.79M \\
Looped TF length-gen (n-RASP-L)~\citep{ax2409.15647} & Fixed-depth weight tying (NTP-Loop) and pause tokens help somewhat but are clearly worse than adaptive depth; pause tokens add parallel not sequential compute. & P, D & small \\
Looped TF length-gen (n-RASP-L)~\citep{ax2409.15647} & Looped transformers with step-dependent final-answer supervision and adaptive loop count length-generalize far beyond training length (parity: trained to 20 digits, near-perfect at 40+). & D, Dat & small GPT-2 blocks (256-dim) \\
Loopie~\citep{ax2607.16051} & Layer-loop (per-layer repetition) beats model-loop (whole-stack repetition) on MoE backbones at the same parameters and FLOPs, but only after about 1.2T tokens; model-loop leads early. & P, IF, Dat & 6B-A0.6B \\
Loopie~\citep{ax2607.16051} & Marginal benefit of extra loops decays fast under fixed pre-training compute; R=2 chosen; the N x layer vs N x loop figure is not compute-matched for N=2. & P, D & 0.25B active \\
LoopLM (Closing the Loop) /\allowbreak{} Looped Qwen~\citep{ax2610.00673} & At matched inference compute (3 loops vs a 72-layer dense twin on the same data), the dense model is ahead or level on every benchmark; the looped model's advantage is parameter count. & IF, D, Dat & 1.4B looped vs 3.9B dense \\
LoopLM (Closing the Loop) /\allowbreak{} Looped Qwen~\citep{ax2610.00673} & Think-at-Hard (2511.08577) could not be reproduced on this training mixture: repeated-step loss never fell below first-step loss. & U & not stated \\
LOTUS~\citep{ax2606.31779} & Trained loop depth matters strongly and extra test-time loops beyond the trained depth do not help. & P, Dat & 3B \\
LOTUS~\citep{ax2606.31779} & At 3B, the looped padded model with parallel gold-CoT supervision brings latent reasoning to within about 1 point of explicit CoT on GSM8K (same backbone, same data), unlike prior latent methods whose gap widens with scale. & P, Dat & 3B \\
LPT /\allowbreak{} H-LPT~\citep{ax2610.02383} & Memory sharing improves quality over a matched-depth looped baseline while reducing KV memory 2.0-4.7x with under 1\% more FLOPs. & P, Dat & 1B \\
LPT /\allowbreak{} H-LPT~\citep{ax2610.02383} & Hybrid H-LPT beats a same-size standard Transformer while using 76-79\% less context memory at five recursions. & P, Dat & 1B \\
MoEUT~\citep{ax2405.16039} & MoEUT slightly outperforms dense Transformers with the same parameter count on C4 perplexity, and the gap grows with scale (up to about 1B). & P, Dat & 1040M \\
MoEUT~\citep{ax2405.16039} & MoEUT slightly outperforms dense Transformers with the same parameter count (same d\_model and number of layers, same training steps), and the gap grows with scale, while using fewer MACs. & P, Dat & 1040M \\
MoR~\citep{ax2507.10524} & At equal training FLOPs, MoR (expert-choice, Nr=2) beats the vanilla Transformer in average few-shot accuracy with about half the parameters, because routing lets it process more tokens in the same budget. & P, TF, Dat & 360M base (167M MoR) \\
MoR~\citep{ax2507.10524} & In isoFLOP sweeps MoR beats fixed-depth recursive at every size but underperforms vanilla at 135M; it matches or exceeds vanilla only above 360M. & TF & 1.7B \\
Nanbeige4.2-3B~\citep{ax2607.22083} & Two passes give the best trade-off, retaining about 75\% of a standard Transformer's token efficiency; more passes add little, slow training and destabilize optimization. & P &  \\
Nanbeige4.2-3B~\citep{ax2607.22083} & Sharing the KV cache across loop passes halves KV memory but consistently yields lower gains than the full non-sharing loop, so it was not used. & P &  \\
\citep{ax2610.02741} & Increasing Ouro loop depth from the default T=4 to T=8 lowers CoT monitorability on Logic, Science and Engineering Cue Answer tasks under stress-test prompts, while other tasks are stable or change differently. & P & Ouro-2.6B-Thinking \\
\citep{ax2610.02741} & The monitorability decline is not explained by changes in verification pass rate across the 32 model x task x stress-test cells. & P & Ouro-2.6B-Thinking \\
Ouro~\citep{ax2510.25741} & On the Mano modular-arithmetic manipulation task, looped models beat iso-parameter non-looped models for all k and often beat the iso-FLOP (12-layer, iso-effective-depth) baseline. & P, D & 12 effective layers, hidden 1024 (synthetic) \\
Ouro~\citep{ax2510.25741} & In the small-scale controlled sweep (20B tokens), a standard transformer whose layer count is multiplied by the recurrent step always beats the LoopLM; the gap grows with loops and shrinks with model size. & D, Dat & 1.36B \\
Parcae~\citep{ax2604.12946} & At fixed parameter count and fixed training FLOPs, the compute-optimal looped frontier (fewer tokens, more recurrence) achieves strictly lower loss and higher Core than fixed-depth models that spend the FLOPs on more data. & P, TF & 370M, up to 1.28e20 FLOPs \\
Parcae~\citep{ax2604.12946} & Compute-optimal training increases mean recurrence and data in tandem following power laws. & P, TF & 370M \\
PLT~\citep{ax2510.24824} & Sequential looping roughly doubles decode latency and KV memory at L=2; CLP keeps accuracy and cuts latency 37\% vs naive loop at bs=4. & P, Dat & 680M activated / 13B total MoE \\
PLT~\citep{ax2510.24824} & Sharing the first loop's KV cache halves KV memory but costs 3.4 average points; gated SWA recovers it with 1.4\% extra KV. & P, Dat & 680M activated / 13B total MoE \\
PonderLM~\citep{ax2505.20674} & PonderPythia trained with the same architecture, tokenizer, data and hyperparameters as official Pythia outperforms it at every size; 2.8B surpasses Pythia-6.9B (but each token costs about s+1 forward passes). & P, Dat & 2.8B, 300B tokens \\
PonderLM~\citep{ax2505.20674} & At equal extra steps, pondering beats a plain looped Transformer, pause tokens and hidden-state feedback; 3 pondering steps beat 6-loop baselines. & P, Dat & 70M, 30B tokens \\
PonderNet~\citep{ax2107.05407} & On parity extrapolation (more ones than in training), PonderNet reaches near-perfect accuracy while ACT stays at chance, and it increases thinking steps to about 5 (roughly twice interpolation). & P & 128-unit RNN \\
PonderNet~\citep{ax2107.05407} & With the same UT transformer architecture, PonderNet solves all 20 bAbI tasks with far fewer steps than UT with ACT (1658 vs 10161). & P & UT small \\
Relaxed Recursive Transformers~\citep{ax2410.20672} & With LoRA rank 512, the relaxed recursive Gemma matches the original full-size Gemma 2B (58.4 vs 58.6) with fewer parameters and 60B uptraining tokens. & D, U & Gemma 2B \\
Relaxed Recursive Transformers~\citep{ax2410.20672} & Continuous depth-wise batching with early exiting gives up to 2.66x simulated generation speedup over vanilla models, but this rests on oracle exits and hypothetical throughput. & D & Gemma 2B -$>$ 1B \\
Retrofitted Recurrence~\citep{ax2511.07384} & A curriculum that slowly raises the mean recurrence keeps loss while reducing training cost; scheduling the max backprop depth instead is FLOP-efficient but worse per step. & TF, Dat & \textasciitilde{}1B, \textasciitilde{}1e20 FLOPs runs \\
Retrofitted Recurrence~\citep{ax2511.07384} & At equal training FLOPs on the same data order, retrofitted recurrent models with fewer unique parameters beat the non-recurrent parent on GSM8K/MATH, and keep improving where the parent plateaus. & TF, Dat & \textasciitilde{}1B, \textasciitilde{}50B tokens \\
RINS~\citep{ax2502.07503} & Stochastic RINS with linear adapters improves the LM even when recursion is disabled at test time (same training compute, parameters and inference cost). & P, TF, IF & 600M \\
RINS~\citep{ax2502.07503} & At matched parameter count and training FLOPs, RINS outperforms all other sharing strategies and longer-context baselines, and its advantage grows with compute. & P, TF & 600M (sweep), 1B \\
Saunshi et al. (k$\otimes$L looping)~\citep{ax2502.17416} & At LM scale, looped models have worse perplexity than the iso-depth baseline but recover far more of the gap on reasoning-heavy downstream groups (open-book QA, math word problems, reasoning primitives) than on memorization (closed-book QA). & P, D, Dat & 24-layer 1B (1.5B in appendix) \\
Saunshi et al. (k$\otimes$L looping)~\citep{ax2502.17416} & On synthetic reasoning tasks, a k-layer block looped L times nearly matches a kL-layer non-looped model and far outperforms the k-layer model. & P, D & 12-layer synthetic models \\
SMELT~\citep{ax2609.01343} & With per-token FLOPs, total parameters and KV cache matched, the looped MoE (SMELT) reaches lower loss than the unlooped MoE at every scale and sparsity level. & P, TF, IF, Dat & 1.6B active / 54B total, \textasciitilde{}215B tokens \\
SMELT~\citep{ax2609.01343} & Downstream gains exceed what the validation-loss gap predicts, and this excess grows with scale; at 1.6B the largest excess shifts to Symbolic Problem Solving. & P, IF, Dat & 1.6B active / 54B total \\
Sparse Layers /\allowbreak{} Looped-MoE~\citep{ax2605.09165} & At matched training compute and effective depth, dense looping scales strictly worse than the dense baseline, while Looped-MoE scales better than the dense baseline. & TF, D & 305M active, 1e18 FLOPs \\
Sparse Layers /\allowbreak{} Looped-MoE~\citep{ax2605.09165} & Looping the same architecture always reduces quality at matched compute: non-looped MoE beats Looped-MoE and Base beats Looped; all four share roughly the same slope and differ in offset. & TF, D & 305M active \\
Takase-Kiyono sharing (SEQUENCE/\allowbreak{}CYCLE/\allowbreak{}CYCLE(REV))~\citep{ax2104.06022} & Universal (one shared layer, widened to match parameters) beats Vanilla in MT but trains more than twice as slow; partial sharing (M=6, N=12) is faster than Universal and slightly better. & P, Dat & 61M \\
Takase-Kiyono sharing (SEQUENCE/\allowbreak{}CYCLE/\allowbreak{}CYCLE(REV))~\citep{ax2104.06022} & In language modeling, full one-layer sharing is worse than no sharing, while partial sharing (N=12, M=6) beats the 6-layer Vanilla at equal parameters; SEQUENCE is best. & P, Dat & 121M (234M for SEQUENCE large) \\
\citep{ax2102.11011} & Recurrent models show the same hierarchical behaviors and depth benefits as feed-forward networks despite reusing filters. & D &  \\
\citep{ax2102.11011} & Recurrent networks can closely emulate non-recurrent deep models, often with far fewer parameters. & D &  \\
TRM~\citep{ax2510.04871} & Backpropagating through a full recursion process instead of HRM's 1-step gradient is the single largest gain (56.5\% to 87.4\% on Sudoku-Extreme). & Dat & 5M \\
TRM~\citep{ax2510.04871} & Fewer layers with proportionally more recursions generalizes better: 2 layers beat 4 layers at roughly matched emulated depth; adding layers overfits. & Dat & 10M vs 5M \\
URM~\citep{ax2512.14693} & At equal FLOPs (32x), a weight-shared UT (4 layers x 8 loops) beats a 32-layer vanilla Transformer on ARC-AGI-1 pass@1 (40.0 vs 23.75); vanilla depth/width scaling saturates or degrades. & IF, D & 32 layers x 512 hidden vanilla \\
URM~\citep{ax2512.14693} & URM reaches 53.8\% pass@1 on ARC-AGI-1 and 16.0\% on ARC-AGI-2, above re-run TRM (40.0) and HRM (34.4). & Dat & 4 layers, hidden 512 \\
UT~\citep{ax1807.03819} & Dynamic halting outperforms fixed-depth UTs even when the fixed models use as many or more steps than the halting model's average (8 and 9 vs 8.2), suggesting halting acts as a regularizer. & D & UT base \\
UT~\citep{ax1807.03819} & UT base beats Transformer base by 0.9 BLEU on WMT14 En-De at approximately the same parameter count (so UT spends more compute per parameter). & P & Transformer-base parameter count (not stated numerically) \\
\citep{ax2609.36636} & With extra test-time loops, looped models trained at the same effective depth and training FLOPs as a 20-layer NonLoop exceed its reasoning score (NonLoop 30.25; 5x4 reaches 31.88 at L(r)=30), using fewer unique layers but more inference compute. & TF, IF, D & Llama3.1-1B config \\
\citep{ax2609.36636} & A non-recurrent coda mitigates knowledge collapse under under-unrolling; prelude-heavy allocation better supports reasoning extrapolation; near the horizon allocation matters little. & P, D & Llama3.1-1B and Qwen3-0.6B configs \\
\end{longtable}
